\section*{Chapitre \ref{ch:b1:titration-methods} --- Méthodes de titrage et courbes de titrage}

\begin{solution}{exo:b1:titration-methods:1}
\ce{H3O+ + OH- -> 2H2O}, $K = 1/K_e = 10^{14.00}$. $C = 12.4 \times
0.100/20.0 = \qty{0.0620}{mol/L}$.
\end{solution}

\begin{solution}{exo:b1:titration-methods:2}
À l'\omterm{def:b1:titration-methods:titration}{équivalence}, la solution est une solution de chlorure d'ammonium à
\qty{0.050}{mol/L}~: $\mathrm{pH} = \frac12(9.25 - \log 0.050) = 5.28$.
Le rouge de méthyle (3.8--5.8) le contient~; l'hélianthine, dont la zone
(2.5--4.5) n'est atteinte qu'après l'\omterm{def:b1:titration-methods:titration}{équivalence}, quand le pH descend,
virerait trop tard.
\end{solution}

\begin{solution}{exo:b1:titration-methods:3}
\ce{MnO4- + 5Fe^2+ + 8H+ -> Mn^2+ + 5Fe^3+ + 4H2O}~;
$C = 5 \times 0.0200 \times 12.6/10.0 = \qty{0.126}{mol/L}$.
\end{solution}

\begin{solution}{exo:b1:titration-methods:4}
$C = 0.0100 \times 14.2/50.0 = \qty{2.84e-3}{mol/L}$ de calcium et de
magnésium réunis. Les deux réagissent avant la fin du \omterm{def:b1:titration-methods:titration}{titrage}, qui ne
donne que leur somme~: c'est un \omterm{def:b1:titration-methods:kinds}{titrage simultané}.
\end{solution}

\begin{solution}{exo:b1:titration-methods:5}
0~mL~: \omterm{def:b1:predominant-reaction:strong-acid}{acide faible}, $\mathrm{pH} = \frac12(4.76 + 1.00) = 2.88$.
5.0~mL~: \omterm{def:b1:titration-methods:titration}{demi-équivalence}, 4.76. 10.0~mL~: 8.73
(\cref{prop:b1:titration-methods:ph-at-equivalence}). 15.0~mL~: excès
d'hydroxyde de \qty{0.50}{mmol} dans \qty{25.0}{mL},
$[\ce{OH-}] = 0.020$, $\mathrm{pH} = 12.30$. Tout concorde avec la courbe
calculée à 0.01 près.
\end{solution}

\begin{solution}{exo:b1:titration-methods:6}
Les $\mathrm{p}K_a$ 2.15 et 7.21 diffèrent de plus de 4, de même que
7.21 et 12.34~: deux \omterm{def:b1:titration-methods:titration}{équivalences}, à 10.0 et à \qty{20.0}{mL}.
Première~: solution de \ce{H2PO4-},
$\mathrm{pH} \approx \frac12(2.15 + 7.21) = 4.68$~; seconde~:
\ce{HPO4^2-}, $\frac12(7.21 + 12.34) = 9.78$ (la courbe exacte donne 4.71
et 9.66, les formules négligeant la dilution et l'eau). La troisième
acidité ($\mathrm{p}K_a$ 12.34) est trop proche de celle de l'eau~:
l'hydroxyde ajouté ne réagit que partiellement et le pH monte sans saut.
\end{solution}

\begin{solution}{exo:b1:titration-methods:7}
Acide chlorhydrique~: $6.2 \times 0.100/10.0 = \qty{0.062}{mol/L}$~;
acide éthanoïque~: $(14.6 - 6.2) \times 0.100/10.0 = \qty{0.084}{mol/L}$.
À mi-chemin, la moitié de l'acide éthanoïque est titrée~:
$\mathrm{pH} = \mathrm{p}K_a = 4.76$.
\end{solution}

\begin{solution}{exo:b1:titration-methods:8}
$V_{\mathrm{eq}} = \qty{10.0}{mL}$. 2.0~mL~: $E = 0.77 + 0.059\log(2/8) =
0.73$~V~; 9.0~mL~: $0.77 + 0.059\log 9 = 0.83$~V~; 11.0~mL~: excès de
permanganate de \qty{0.020}{mmol}, \ce{Mn^2+} \qty{0.200}{mmol}, $E = 1.51 +
\frac{0.059}{5}\log 0.10 = 1.50$~V~; 20.0~mL~: 1.51~V.
\end{solution}

\begin{solution}{exo:b1:titration-methods:9}
0~mL~: $\kappa = (349.8 + 76.2) \times 0.0100 = \qty{4.26}{mS/cm}$.
10.0~mL~: \ce{Na+} et \ce{Cl-} à $1.00/110 = \qty{9.09e-3}{mol/L}$,
$\kappa = (50.3 + 76.2) \times \num{9.09e-3} = 1.15$. 20.0~mL~: \ce{Na+}
$0.0167$, \ce{Cl-} et \ce{OH-} $0.00833$~mol/L, $\kappa = 0.84 + 0.635 + 1.65
= 3.13$~mS/cm. Pentes (dilution négligée, \qty{1}{mL} ajoute
\qty{1.0e-3}{mol/L})~: $(50.3 - 349.8) \times 10^{-3} = -0.30$ et
$(50.3 + 198.3) \times 10^{-3} = +0.25$~mS/cm par mL.
\end{solution}

\begin{solution}{exo:b1:titration-methods:10}
En $V_{\mathrm{eq}} - v$, l'excès d'acide vaut $C_Bv$~; en
$V_{\mathrm{eq}} + v$, l'excès de base vaut $C_Bv$~; dans le même
volume, $[\ce{H3O+}]$ avant est égal à
$[\ce{OH-}]$ après, donc $\mathrm{pH}(V_{\mathrm{eq}} - v) +
\mathrm{pH}(V_{\mathrm{eq}} + v) = 14$~: symétrie par rapport au point
(V$_{\mathrm{eq}}$,
7). Saut~: excès de \qty{1.0e-5}{mol} dans \qty{20}{mL}, pH de 3.30 à
10.70, soit 7.4 unités. Concentrations divisées par 100~:
$\num{1.0e-7}$~mol dans \qty{20}{mL}, pH de 5.3 à 8.7, soit 3.4
unités~: une fin de \omterm{def:b1:titration-methods:titration}{titrage} bien moins nette.
\end{solution}

\begin{solution}{exo:b1:titration-methods:11}
Avant~: \ce{Fe^3+} formé $= 5C_{\mathrm{Mn}}V$, \ce{Fe^2+} restant $=
5C_{\mathrm{Mn}}(V_{\mathrm{eq}} - V)$, donc $E = 0.77 + 0.059\log\frac{V}{V_{\mathrm{eq}}
- V}$~: en $V_{\mathrm{eq}}/2$, $E = 0.77$. Après~: excès de \ce{MnO4-} $\propto
V - V_{\mathrm{eq}}$, \ce{Mn^2+} $\propto V_{\mathrm{eq}}$~: $E = 1.51 +
\frac{0.059}{5}\log\frac{V - V_{\mathrm{eq}}}{V_{\mathrm{eq}}}$ (pH 0)~;
en $2V_{\mathrm{eq}}$, 1.51~V. À l'\omterm{def:b1:titration-methods:titration}{équivalence},
$(0.77 + 5 \times 1.51)/6 =
1.39$~V. À pH 1, le couple du permanganate est plus bas de
$0.059 \times 8/5 =
0.094$~V~: $E_{\mathrm{eq}} = (0.77 + 5 \times 1.416)/6 = 1.31$~V.
\end{solution}

\begin{solution}{exo:b1:titration-methods:12}
\ce{NH4+ + OH- -> NH3 + H2O} a pour constante $K = 10^{14.00 - 9.25} = 10^{4.75}$~: la
réaction n'est pas complète près de l'\omterm{def:b1:titration-methods:titration}{équivalence} et le saut est faible,
vers pH 11 (pH à l'\omterm{def:b1:titration-methods:titration}{équivalence} $14 - \frac12(4.75 + 1.30) = 11.0$).
\omterm{def:b1:titration-methods:kinds}{Titrage en retour}~: ajouter $n_1$ d'hydroxyde de sodium (excès), faire
bouillir pour chasser l'ammoniac formé, puis titrer l'hydroxyde restant
par l'acide chlorhydrique ($n_2$, saut net acide fort--base forte)~:
$n(\ce{NH4+}) = n_1 - n_2$.
\end{solution}

\begin{solution}{pb:b1:titration-methods:1}
\textbf{1.} \ce{C6H6O6 + 2H+ + 2e- -> C6H8O6}.
\textbf{2.} \ce{C6H8O6 + I2 -> C6H6O6 + 2I- + 2H+}.
\textbf{3.} 0 et $-$I.
\textbf{4.} \ce{I2 + 2S2O3^2- -> 2I- + S4O6^2-}, $\log K = 2(0.53 - 0.02)/0.059
= 17$.
\textbf{5.} La couleur bleue exige du diiode~; elle disparaît avec le
dernier diiode, à l'\omterm{def:b1:titration-methods:titration}{équivalence}.
\textbf{6.} Des pertes à l'air diminueraient le résultat~; ajouté d'un
coup et en excès, le diiode oxyde l'acide ascorbique rapidement et
complètement.
\textbf{7.} Une réaction rapide et \omterm{def:b1:extent-q-and-k:quantitative}{quantitative} pendant tout le
\omterm{def:b1:titration-methods:titration}{titrage}, sans oxydation par l'air pendant l'ajout lent du \omterm{def:b1:titration-methods:titration}{réactif
titrant}, et une fin de \omterm{def:b1:titration-methods:titration}{titrage} au premier excès de diiode.
\textbf{8.} Excès de diiode $n_R$~; réaction avec l'acide ascorbique~;
\omterm{def:b1:titration-methods:titration}{titrage} du diiode restant par le thiosulfate.
\textbf{9.} Sinon, il resterait de l'acide ascorbique et l'excès de
diiode, nul, ne dirait pas combien. Ici, \qty{2.775e-4}{mol} ont réagi
sur \qty{5.000e-4}{mol}~: il y a bien un excès.
\textbf{10.} Les solutions de thiosulfate évoluent lentement dans le
temps~; sa concentration intervient directement dans le résultat.
\textbf{11.} 10.
\textbf{12.} Une pipette jaugée de \qty{20.00}{mL}.
\textbf{13.} $20.00 \times 10^{-3} \times 0.0250 = \qty{5.000e-4}{mol}$.
\textbf{14.} $8.90 \times 10^{-3} \times 0.0500 = \qty{4.450e-4}{mol}$.
\textbf{15.} La moitié~: \qty{2.225e-4}{mol}.
\textbf{16.} $5.000 - 2.225 = \qty{2.775e-4}{mol}$.
\textbf{17.} Mole à mole~: \qty{2.775e-4}{mol}.
\textbf{18.} $\times 10$~: \qty{2.775e-3}{mol}, $\times 176.0$~:
\qty{0.488}{g}.
\textbf{19.} Environ \qty{2}{\percent} en dessous de l'étiquette.
\textbf{20.} $n = \bigl(C_IV_I - \frac12C_TV_T\bigr)\,V_{\mathrm{fiole}}/V_{\mathrm{aliquote}}$.
\textbf{21.} Diiode introduit~: pipette $0.03/20.00 = \qty{0.15}{\percent}$
et concentration \qty{0.2}{\percent}, environ \qty{0.25}{\percent},
soit \qty{1.3e-6}{mol}. Excès~: burette $0.05/8.90 = \qty{0.56}{\percent}$
et concentration \qty{0.2}{\percent}, environ \qty{0.60}{\percent}, soit
\qty{1.3e-6}{mol}.
\textbf{22.} Les incertitudes absolues se combinent, alors que la
différence est plus petite que chacune des deux quantités~:
$\sqrt{1.3^2 + 1.3^2} \times 10^{-6} =
\qty{1.8e-6}{mol}$, soit \qty{0.66}{\percent} de \qty{2.775e-4}{mol},
plus que chacun des deux termes.
\textbf{23.} Avec la fiole (\qty{0.1}{\percent}) et la pipette de la
partie aliquote (\qty{0.2}{\percent})~:
$\sqrt{0.66^2 + 0.1^2 + 0.2^2} = \qty{0.70}{\percent}$.
\textbf{24.} $0.488 \pm 0.003$~g~: \textbf{\qty{0.49}{g} d'acide
ascorbique par comprimé}.
\end{solution}
